\section{Density martingales and infinite product measures}
\label{sec:annex-p15-density-martingales}

\begin{annexitem}{annex:p15.density-martingale}{thm:p15-density-martingale}{Restriction identities under the correct probability}
\begin{proof}
For \(A\in\mathcal F_n\), restriction of the integration space gives
\[
                 \int_A X_n\,d\nu
                  =\int_A X_n\,d\nu_n=\mu_n(A)=\mu(A).
\]
In particular \(X_n\) is integrable under \(\nu\) with integral \(c\).
If \(m\geq n\), the same identity applied at time \(m\) gives
\(\int_A X_m\,d\nu=\int_A X_n\,d\nu\) for all
\(A\in\mathcal F_n\). This is exactly
\(E_\nu(X_m\mid\mathcal F_n)=X_n\), by the defining uniqueness of
conditional expectation. Adaptedness is part of the derivative choice.
The nonnegative martingale convergence theorem gives finite convergence
\(\nu\)-almost surely. Fatou's lemma gives
\(E_\nu X\leq\liminf_nE_\nu X_n=c\).
All assertions concern \(\nu\); this argument makes no claim that the
limit is finite under \(\mu\).
\end{proof}
\end{annexitem}

\begin{annexitem}{annex:p15.density-mixture-reference}{thm:p15-density-mixture-reference}{The bounded mixture construction and its measure transfer}
\begin{proof}
Since \(c\) is finite and nonnegative, \(\rho\) is a probability.
Both \(\mu\) and \(\nu\) are absolutely continuous with respect to
\(\rho\). Their restricted densities exist by
\cref{thm:p15-radon-nikodym}. The inequalities
\(\mu_n\leq(c+1)\rho_n\), \(\nu_n\leq(c+1)\rho_n\) imply
\(0\leq Y_n,Z_n\leq c+1\), \(\rho\)-almost surely:
a set where a density exceeds its bound would contradict that measure
inequality by integration. The sum identity in
\cref{thm:p15-density-calculus} gives \(Y_n+Z_n=c+1\).
Apply \cref{thm:p15-density-martingale} with base \(\rho\), separately
to the two finite measures. Bounded convergence gives their
almost-sure and \(L^1(\rho)\) limits \(Y,Z\), retaining the sum identity.

For \(A\in\mathcal F_m\) and \(n\geq m\),
\(\int_A Y_n\,d\rho=\mu(A)\). Passing to the limit gives
\(\int_A Y\,d\rho=\mu(A)\). The union
\(\mathcal G=\bigcup_m\mathcal F_m\) is an algebra: finitely many of
its members belong to a common \(\mathcal F_m\).
It contains \(\Omega\) and is closed under finite intersections.
The two finite measures \(A\mapsto\int_A Y\,d\rho\) and \(\mu\)
therefore agree on the generating pi-system \(\mathcal G\);
the pi-lambda uniqueness theorem extends equality to \(\mathcal F\).
The identical argument with \(Z_n\) gives \(\nu=Z\rho\).

Local absolute continuity also gives \(\rho_n\ll\nu_n\):
a \(\nu_n\)-null event has both \(\mu_n\) and \(\nu_n\) mass zero.
On the event \(\{Z_n=0\}\in\mathcal F_n\),
\(\nu_n(\{Z_n=0\})=\int_{\{Z_n=0\}}Z_n\,d\rho_n=0\);
hence \(\rho_n(\{Z_n=0\})=0\).
The density-chain formula gives
\(Y_n=X_nZ_n\), \(\rho_n\)-almost surely, so \(X_n=Y_n/Z_n\)
there. Countably many exceptional events have \(\rho\)-measure zero.

Take an outcome outside those events where both bounded martingales
converge and \(Y+Z=c+1\).
If \(Z>0\), the ratio converges to \(Y/Z\).
If \(Z=0\), then \(Y=c+1>0\), while the positive \(Z_n\) tend to
zero. Thus \(X_n=Y_n/Z_n\) tends to positive infinity.
This proves the full sequence assertions, and hence the stated limsup
identities. It also covers \(c=0\): no normalization by a zero mass
was made. Finite measurable choices off the exceptional set do not
change the measure identities.
\end{proof}
\end{annexitem}

\begin{annexitem}{annex:p15.density-lebesgue-decomposition}{thm:p15-density-lebesgue-decomposition}{The exact regular and singular terms}
\begin{proof}
Use \(Y,Z,\rho\) from the mixture lemma, and define a measurable function
\(W\) by \(W=1/Z\) where \(Z>0\) and \(W=0\) where \(Z=0\).
This piecewise definition never evaluates \(1/0\).
Since \(1=ZW+\mathbf1_{\{Z=0\}}\),
\[
 \mu(A)=\int_A Y\,d\rho
       =\int_A YWZ\,d\rho+\int_A\mathbf1_{\{Z=0\}}Y\,d\rho.
\]
All terms are nonnegative. Change of measure using \(\nu=Z\rho\)
makes the first term \(\int_A YW\,d\nu\).
Now \(\nu(Z=0)=\int_{\{Z=0\}}Z\,d\rho=0\), and on \(Z>0\)
the mixture lemma gives \(YW=X\), \(\rho\)-almost surely.
Since \(\nu\ll\rho\), this also holds \(\nu\)-almost surely.
The first term is therefore \(\int_A X\,d\nu\).

The second term is \(\mu(A\cap\{Z=0\})\).
The mixture lemma identifies \(\{Z=0\}\) and \(\{X=\infty\}\)
up to a \(\rho\)-null event, hence also up to a \(\mu\)-null event.
This proves the displayed decomposition on every \(A\in\mathcal F\).
The integral term is a measure absolutely continuous with respect to
\(\nu\), by nonnegative integration. The other term is concentrated on
\(\{X=\infty\}\), which is \(\nu\)-null by the density martingale lemma.
It is singular with respect to \(\nu\).
Uniqueness is exactly \cref{thm:p15-lebesgue-decomposition}, and density
uniqueness identifies \(X=d\mu_r/d\nu\).
The proof never transfers a merely \(\nu\)-almost-sure statement to
\(\mu\) without the common reference.
\end{proof}
\end{annexitem}

\begin{annexitem}{annex:p15.density-uniform-integrability}{thm:p15-density-uniform-integrability}{No lost mass exactly characterizes uniform integrability}
\begin{proof}
Taking \(A=\Omega\) in the decomposition gives
\(\mu_s(\Omega)=c-E_\nu X\).
A nonnegative measure is zero exactly when its total mass is zero.
Thus preservation of mass is equivalent to \(\mu_s=0\), and this is
equivalent to \(\mu\ll\nu\) by uniqueness of the decomposition.

If \(\mu\ll\nu\), put \(f=d\mu/d\nu\in L^1(\nu)\).
For \(A\in\mathcal F_n\) the density identity says
\(\int_A f\,d\nu=\int_A X_n\,d\nu\).
Hence \(X_n=E_\nu(f\mid\mathcal F_n)\).
The family of conditional expectations of one integrable variable is
uniformly integrable by \cref{thm:p15-ui-conditional-family}.
The generated sigma-field is all of \(\mathcal F\), so
\cref{thm:p15-levy-upward} also gives \(X_n\to f\) almost surely
and in \(L^1\). Almost-sure uniqueness identifies \(f=X\).
Alternatively uniform integrability and the already proved convergence
give \(L^1\) convergence by \cref{thm:p15-vitali-convergence}.
Conversely \(L^1\) convergence gives
\(|E_\nu X-c|\leq E_\nu|X-X_n|\to0\).
This proves every equivalence and the conditional-expectation identity.

If \(\mu\perp\nu\), uniqueness of the Lebesgue decomposition forces
\(\mu_r=0\), so \(X=0\), \(\nu\)-almost surely.
If \(X=0\), the integral component is zero and \(\mu=\mu_s\perp\nu\).
Finally, when \(\mu=X\nu\), strict positivity of \(X\)
\(\nu\)-almost surely implies \(\nu\ll\mu\): a zero integral of \(X\)
on an event implies that event is \(\nu\)-null.
Conversely \(\mu\sim\nu\) makes \(\{X=0\}\) a \(\nu\)-null event
because its \(\mu\)-mass is zero. Mutual absolute continuity also
implies \(\mu\ll\nu\), hence preservation of mass.
\end{proof}
\end{annexitem}

\begin{annexitem}{annex:p15.partition-density-approximation}{thm:p15-partition-density-approximation}{Finite cells and the missing singular mass}
\begin{proof}
Every event in the sigma-field of a finite partition is a union of its
cells. If a cell has zero \(\nu\)-mass but positive \(\mu\)-mass,
local absolute continuity fails. If every \(\nu\)-null cell is also
\(\mu\)-null, the stated ratio function integrates on each cell to
its \(\mu\)-mass, including zero on null cells.
Summing over any cell union proves both absolute continuity and the
density formula. Refinement gives the increasing filtration required
for the preceding results.

For dyadic cells each denominator is \(2^{-n}>0\).
The numerator is \((1-\theta)2^{-n}\) plus \(\theta\) in the cell
containing \(x_0\), proving the explicit formula.
The cells \(I_n(x_0)\) are nested, contain \(x_0\), and have diameters
tending to zero. Their intersection is \(\{x_0\}\): any distinct point
is eventually at distance greater than the cell diameter.
Thus every other point eventually leaves them. The two limits follow.
Integrating the formula gives
\(1-\theta+\theta2^n2^{-n}=1\), whereas the limiting density has
\(\nu\)-integral \(1-\theta\), since a singleton is Lebesgue-null.
The singular part is exactly \(\theta\delta_{x_0}\).

These dyadic partitions generate the relative Borel sigma-field.
Their cells are Borel. Conversely, for every point of a relatively
open set, a sufficiently small containing dyadic cell lies in that set.
The open set is therefore a union from the countable collection of
dyadic cells, so belongs to their generated sigma-field.

For the final example the trivial sigma-field sees only total masses,
both equal to one. Therefore \(X_n=1\) is the derivative at every
time and \(X=1\). But on the Borel singleton \(A=\{x_0\}\),
the proposed full decomposition with this \(X\) would give
\(\int_A1\,d\nu+\mu(A\cap\varnothing)=0\), while \(\mu(A)=1\).
The local formulas recover only the restrictions to their generated
trivial sigma-field, not the full Borel law.
\end{proof}
\end{annexitem}

\begin{annexitem}{annex:p15.product-density-filtration}{thm:p15-product-density-filtration}{Cylinder integration and exact square-root distances}
\begin{proof}
For a cylinder \(A=\{\xi_1\in B_1,\ldots,\xi_n\in B_n\}\),
independence and nonnegative product integration give
\[
 \int_A\prod_{m=1}^nq_m(\xi_m)\,d\nu
       =\prod_{m=1}^n\int_{B_m}q_m\,d\beta_m
       =\prod_{m=1}^n\alpha_m(B_m)=\mu(A).
\]
The total integral is one. The two finite measures therefore agree
on the cylinder pi-system generating \(\mathcal F_n\), and hence on
all of \(\mathcal F_n\). This proves the density and local absolute
continuity claims. The full product sigma-field is generated by the
union of these finite-coordinate sigma-fields.

Cauchy--Schwarz gives
\(h_m^2\leq(\int q_m\,d\beta_m)(\int1\,d\beta_m)=1\).
It is positive because a zero integral of \(\sqrt{q_m}\) would give
\(q_m=0\), \(\beta_m\)-almost surely, contradicting
\(\int q_m\,d\beta_m=1\).
Also
\[
                  \int(\sqrt{q_m}-1)^2\,d\beta_m=2(1-h_m).
\]
Consequently \(h_m=1\) exactly when \(q_m=1\) almost surely,
which is equivalent to equality of the two coordinate laws.
Each factor belongs to \((0,1]\), proving existence of \(H_\infty\).

Independence gives \(E_\nu\sqrt{X_n}=\prod_{m\leq n}h_m=H_n\)
and \(E_\nu V_n^2=E_\nu X_n=1\).
For \(l>n\),
\[
 V_lV_n=X_n\prod_{m=n+1}^l\sqrt{q_m(\xi_m)}.
\]
The two factors use disjoint independent coordinate groups, are
nonnegative, and have the finite indicated integrals. Therefore
\(E_\nu(V_lV_n)=\prod_{m=n+1}^lh_m\).
Expanding the square using \(E_\nu V_l^2=E_\nu V_n^2=1\) proves
the exact distance formula. The calculation includes \(n=0\).
\end{proof}
\end{annexitem}

\begin{annexitem}{annex:p15.product-positive-tail}{thm:p15-product-positive-tail}{Both probabilities in the source tail argument}
\begin{proof}
Strict positivity \(\beta_m\)-almost surely and
\(\alpha_m=q_m\beta_m\) imply \(\beta_m\ll\alpha_m\).
Replace \(q_m\) by one on its \(\beta_m\)-null set of zero or nonfinite
values. That set is also \(\alpha_m\)-null.
Now \(q_m\) is finite and positive everywhere, and
\(d\beta_m/d\alpha_m=1/q_m\), since for every Borel \(B\),
\(\int_B(1/q_m)\,d\alpha_m=\int_B1\,d\beta_m\).
The finite-product density proof applied in the reverse direction
makes \(R_n=1/X_n\) a nonnegative martingale under \(\mu\).

For a real sequence \(z_m=\log q_m(\xi_m)\), let \(C\) be the event
that \(\sum_{m\geq1}z_m\) converges to a finite real number.
This is measurable by the countable Cauchy criterion.
Convergence is unchanged by deleting finitely many real summands,
so \(C\in\sigma(\xi_m:m>k)\) for every \(k\).
It is a tail event under either independent product measure, and
Kolmogorov's theorem \cref{thm:p15-kolmogorov-tail-zero-one} gives
\(\nu(C),\mu(C)\in\{0,1\}\).

Under \(\nu\), the positive \(X_n\) converge to a finite nonnegative
\(X\) almost surely. On that event, \(X>0\) exactly when
\(\log X_n=\sum_{m\leq n}z_m\) has a finite limit, by continuity
of logarithm and exponential. Thus \(\nu(X=0)=1-\nu(C)\)
is zero or one.
Under \(\mu\), the reciprocal martingale has a finite nonnegative
limit \(R\) almost surely. If \(R>0\), \(X_n\to1/R\);
if \(R=0\), \(X_n\to\infty\).
The same convergence criterion gives
\(\mu(X=\infty)=\mu(R=0)=1-\mu(C)\in\{0,1\}\).

The density decomposition gives
\(\mu_s(\Omega)=\mu(X=\infty)\).
If this is one, \(\mu=\mu_s\perp\nu\).
If it is zero, \(\mu\ll\nu\) and \(E_\nu X=1\).
Then \(\nu(X=0)\ne1\); its zero-one value is therefore zero,
so \(X>0\), \(\nu\)-almost surely. The preceding density criterion
gives \(\nu\ll\mu\).
The alternatives are exclusive for probability measures: equivalent
nonzero measures cannot also be concentrated on disjoint null supports.
This completes the tail route without inferring preservation of mass
from strict positivity alone.
\end{proof}
\end{annexitem}

\begin{annexitem}{annex:p15.product-affinity-criterion}{thm:p15-product-affinity-criterion}{The source's square-root estimate, including zero densities}
\begin{proof}
The density martingale theorem gives \(X_n\to X<\infty\),
\(\nu\)-almost surely. If \(H_\infty=0\), Fatou gives
\[
                  E_\nu\sqrt X
                     \leq\liminf_nE_\nu\sqrt{X_n}
                     =\lim_nH_n=0.
\]
Hence \(X=0\), \(\nu\)-almost surely, and the density decomposition
implies \(\mu\perp\nu\).
It also gives \(\mu(X=\infty)=1\). More precisely the mixture-reference
proof shows actual convergence \(X_n\to\infty\) on the set \(Z=0\)
that carries this singular measure; a limsup alone is not being
upgraded to convergence.

Suppose now \(H_\infty>0\). For \(l>n\), write \(V_j=\sqrt{X_j}\).
Cauchy--Schwarz and \((u+v)^2\leq2u^2+2v^2\) give
\[
 \begin{aligned}
 E_\nu|X_l-X_n|
 &=E_\nu\bigl(|V_l-V_n|(V_l+V_n)\bigr)\\
 &\leq\bigl(E_\nu(V_l-V_n)^2\bigr)^{1/2}
        \bigl(E_\nu(V_l+V_n)^2\bigr)^{1/2}\\
 &\leq2\sqrt{2\left(1-\frac{H_l}{H_n}\right)}.
 \end{aligned}
\]
All denominators \(H_n\) are positive finite products.
Fatou applied with fixed \(n\) and \(l\to\infty\) yields
\[
             E_\nu|X-X_n|
                  \leq2\sqrt{2\left(1-\frac{H_\infty}{H_n}\right)}
                  \longrightarrow0.
\]
Thus the limit has integral one. The density uniform-integrability
criterion gives \(\mu=X\nu\). This argument allows \(q_m=0\);
no logarithms or reciprocal products have been used.

Since \(H_\infty\) is either positive or zero and the probability
measures cannot be both absolutely continuous and singular in the same
direction, the two implications prove both equivalences.
If the coordinate laws are equivalent, then
\[
 \int\sqrt{d\beta_m/d\alpha_m}\,d\alpha_m
       =\int\frac{1}{\sqrt{q_m}}q_m\,d\beta_m=h_m,
\]
using positive versions. Apply the proved criterion with the measures
reversed to obtain \(\nu\ll\mu\) when \(H_\infty>0\).
The one-sided hypothesis alone supplies no such reverse implication.
\end{proof}
\end{annexitem}

\begin{annexitem}{annex:p15.product-affinity-series}{thm:p15-product-affinity-series}{From products to sums}
\begin{proof}
For \(0<x\leq1\),
\[
             1-x\leq-\log x=\int_x^1\frac{dt}{t}.
\]
If \(1/2\leq x\leq1\), the integral is at most \(2(1-x)\).
A positive product limit is equivalent to boundedness of
\(-\log H_n=\sum_{m\leq n}(-\log h_m)\), hence to convergence of
the nonnegative logarithmic series.
Its convergence implies \(\sum_m(1-h_m)<\infty\) by the first bound.
Conversely that latter summability forces \(h_m\to1\), so all but
finitely many \(h_m\) lie in \([1/2,1]\); the second bound proves
convergence of the logarithmic tail. Each omitted initial logarithm
is finite because every \(h_m>0\). Finally the exact square identity
from the finite-affinity lemma gives the second series as twice
the first. The product criterion supplies the claimed measure relation.
\end{proof}
\end{annexitem}

\begin{annexitem}{annex:p15.product-gaussian-comparison}{thm:p15-product-gaussian-comparison}{Completing the square for each normal coordinate}
\begin{proof}
For one coordinate suppress the subscript and let \(f,g\) be the
strictly positive Lebesgue densities of
\(\operatorname{normal}(a,\sigma^2)\) and
\(\operatorname{normal}(b,\tau^2)\).
Their ratio represents the derivative relative to the second law,
so \(h=\int_{\mathbb R}\sqrt{fg}\,dx\).
Set \(D=\sigma^2+\tau^2\) and
\(v=(a\tau^2+b\sigma^2)/D\).
Expanding both sides verifies
\[
 \frac{(x-a)^2}{4\sigma^2}+\frac{(x-b)^2}{4\tau^2}
 =\frac{D}{4\sigma^2\tau^2}(x-v)^2+\frac{(a-b)^2}{4D}.
\]
The prefactor of \(\sqrt{fg}\) is
\((2\pi\sigma\tau)^{-1/2}\).
The already normalized Gaussian integral gives
\[
 \int_{\mathbb R}
       \exp\!\left(-\frac{D}{4\sigma^2\tau^2}(x-v)^2\right)dx
                =2\sigma\tau\sqrt{\frac{\pi}{D}}.
\]
Multiplying the factors gives exactly the displayed \(h\).

Every coordinate pair is equivalent: both laws have strictly positive
Lebesgue densities, so each assigns positive mass to exactly the Borel
sets of positive Lebesgue measure, and \(q_m=f_m/g_m\) is a finite,
strictly positive version.
Suppose \(\sum_m(1-h_m)<\infty\). By
\cref{thm:p15-product-affinity-series}, \(\mu\ll\nu\), and \(H_\infty>0\)
by \cref{thm:p15-product-affinity-criterion}. With the measures
reversed, the derivative of \(\beta_m\) relative to \(\alpha_m\) is
\(1/q_m\), and the reverse-measure computation in the proof of
\cref{thm:p15-product-affinity-criterion} shows its affinity is again
\(h_m\). The reversed square-root product is therefore the same
\(H_\infty>0\), and the criterion applied to the reversed pair gives
\(\nu\ll\mu\). Hence \(\mu\sim\nu\).

Suppose instead \(\sum_m(1-h_m)=\infty\). Every \(h_m\) lies in
\((0,1]\) by \cref{thm:p15-product-density-filtration}, and
\(1-x\leq-\log x\) for \(0<x\leq1\), as in the proof of
\cref{thm:p15-product-affinity-series}. Therefore
\[
 -\log H_n=\sum_{m=1}^n(-\log h_m)
          \geq\sum_{m=1}^n(1-h_m)\longrightarrow\infty,
\]
so \(H_n\to0\) and \(H_\infty=0\). The second line of
\cref{thm:p15-product-affinity-criterion} now gives \(\mu\perp\nu\):
its proof shows \(X_n\to X=0\), \(\nu\)-almost surely, and the density
decomposition then puts all of \(\mu\) on the \(\nu\)-null set
\(\{X=\infty\}\). This proves the second alternative of the first
assertion. The two alternatives exclude each other, because two
equivalent probability measures cannot be carried by disjoint sets one
of which is null for each.

When \(\tau_m=\sigma_m\), the prefactor is one and
\(h_m=\exp(-(a_m-b_m)^2/(8\sigma_m^2))\).
Its finite product is
\(\exp(-\sum_{m\leq n}(a_m-b_m)^2/(8\sigma_m^2))\).
The limit is positive exactly when that nonnegative series is finite.
If the series is finite, \(H_\infty>0\), and the equivalence argument of
the preceding paragraphs applies, since \(\sum_m(1-h_m)<\infty\) by
\cref{thm:p15-product-affinity-series}. If the series diverges,
\(H_\infty=0\), so \(\mu\perp\nu\) by
\cref{thm:p15-product-affinity-criterion}, and singular probability
measures are not equivalent. This proves the simplified criterion in
both directions.

Finite-dimensional normal densities are strictly positive in every
case, including the divergent-series case. For instance, take
\(a_m=1\), \(b_m=0\) and \(\sigma_m=\tau_m=1\) for every \(m\). Then
\(h_m=e^{-1/8}\) and \(H_n=e^{-n/8}\to0\), so the sequence laws are
singular by the case just proved, while for every \(n\) the laws
\(\mu_n\) and \(\nu_n\) of \((\xi_1,\ldots,\xi_n)\) have strictly
positive Lebesgue densities on \(\mathbb R^n\) and are equivalent.
Therefore equality of finite-dimensional supports does not force full
equivalence.
\end{proof}
\end{annexitem}

\begin{annexitem}{annex:p15.product-bernoulli-boundaries}{thm:p15-product-bernoulli-boundaries}{Two-point affinities and the positivity hypothesis}
\begin{proof}
The Bernoulli measures give masses \(p,1-p\) and \(q,1-q\) to
\(1,0\). Their derivative values are \(p/q\) and
\((1-p)/(1-q)\). Integrating their square roots against the second
law gives \(\sqrt{pq}+\sqrt{(1-p)(1-q)}\).
With \(q=1/2\), \(p=1/2+\varepsilon\), expand squares to get
\[
 \begin{aligned}
 2(1-h)
   &=(\sqrt p-\sqrt q)^2
        +(\sqrt{1-p}-\sqrt{1-q})^2\\
   &=\varepsilon^2\left(
       \frac{1}{(\sqrt p+\sqrt q)^2}
       +\frac{1}{(\sqrt{1-p}+\sqrt{1-q})^2}\right).
 \end{aligned}
\]
Every probability appearing under these roots lies in \([1/4,3/4]\).
Each denominator therefore lies in \([1,3]\), so the parenthesized
sum lies in \([2/3,2]\).
Division by two proves the stated bounds.

Since \(|\varepsilon_m|\leq1/4\), every \(p_m\) lies in
\([1/4,3/4]\), and \(q_m=1/2\). Two Bernoulli laws whose success
probabilities lie in \((0,1)\) are equivalent: each gives positive mass
to both points \(0\) and \(1\) and no mass to the rest of the line, so
a Borel set is null for either law exactly when it contains neither
point. Their derivatives are the finite positive values computed above.
By the two-sided bounds, \(\sum_m(1-h_m)<\infty\) exactly when
\(\sum_m\varepsilon_m^2<\infty\).

If \(\sum_m\varepsilon_m^2<\infty\), then \(\sum_m(1-h_m)<\infty\), and
\cref{thm:p15-product-affinity-series} gives \(\mu\ll\nu\) with
\(H_\infty>0\). The coordinate laws being equivalent, the
reverse-measure part of \cref{thm:p15-product-affinity-criterion}
applies: the derivative of the second Bernoulli law relative to the
first takes the values \(q/p\) and \((1-q)/(1-p)\), and integrating its
square roots against the first law gives
\(p\sqrt{q/p}+(1-p)\sqrt{(1-q)/(1-p)}=\sqrt{pq}+\sqrt{(1-p)(1-q)}=h\).
The reversed square-root product is again \(H_\infty>0\), so
\(\nu\ll\mu\), and \(\mu\sim\nu\).

If \(\sum_m\varepsilon_m^2=\infty\), then \(\sum_m(1-h_m)=\infty\).
Since \(0<h_m\leq1\) and \(1-x\leq-\log x\) on \((0,1]\),
\[
 -\log H_n=\sum_{m=1}^n(-\log h_m)\geq\sum_{m=1}^n(1-h_m)
          \longrightarrow\infty,
\]
so \(H_\infty=0\), and \cref{thm:p15-product-affinity-criterion} gives
\(\mu\perp\nu\). Singular probability measures are not equivalent. This
proves that the laws are equivalent exactly when
\(\sum_m\varepsilon_m^2<\infty\).

For \(\varepsilon_m=1/(4m)\), \(|\varepsilon_m|\leq1/4\) and
\(\sum_m\varepsilon_m^2=\sum_m1/(16m^2)<\infty\), so the laws are
equivalent. For \(\varepsilon_m=1/(4\sqrt m)\), again
\(|\varepsilon_m|\leq1/4\), while
\(\sum_m\varepsilon_m^2=\sum_m1/(16m)=\infty\); hence
\(\sum_m(1-h_m)\geq\sum_m\varepsilon_m^2/3=\infty\), \(H_\infty=0\) and
\(\mu\perp\nu\). In both cases \(\varepsilon_m\to0\).

For the final example the first density relative to a fair coordinate
is two at zero and zero at one, and every later density is one.
Thus the likelihood products equal
\(2\mathbf1_{\{\xi_1=0\}}\) for all \(n\geq1\).
The first affinity is \(1/\sqrt2\), and all later affinities are one.
The product criterion gives \(\mu\ll\nu\).
But \(\{\xi_1=1\}\) has \(\mu\)-mass zero and \(\nu\)-mass \(1/2\);
hence \(\nu\not\ll\mu\). It is also exactly \(\{X=0\}\).
That event depends on the first coordinate, so it is not the
tail-convergence event used in the strict-positivity theorem.
\end{proof}
\end{annexitem}

